\documentclass[10pt, letterpaper]{article}


\title{Alex Dils Resume}
% Packages:
\usepackage[
    ignoreheadfoot,
    top=1.05cm,
    bottom=1.05cm,
    left=1.40cm,
    right=1.40cm,
    footskip=0.45cm,
]{geometry}
\usepackage{titlesec}
\usepackage{tabularx}
\usepackage{array}
\usepackage[dvipsnames]{xcolor}
\definecolor{primaryColor}{RGB}{0, 0, 0}
\usepackage{enumitem}
\usepackage{fontawesome}
\usepackage{amsmath}
\usepackage[
    pdftitle={Alex Dils's CV},
    pdfauthor={Alex Dils},
    pdfcreator={LaTeX with RenderCV},
    colorlinks=true,
    urlcolor=primaryColor
]{hyperref}
\usepackage[pscoord]{eso-pic}
\usepackage{calc}
\usepackage{bookmark}
\usepackage{lastpage}
\usepackage{changepage}
\usepackage{paracol}
\usepackage{ifthen}
\usepackage{needspace}
\usepackage{iftex}
\usepackage{etoolbox}

\ifPDFTeX
    \input{glyphtounicode}
    \pdfgentounicode=1
    \usepackage[T1]{fontenc}
    \usepackage[utf8]{inputenc}
    \usepackage{lmodern}
    \usepackage{charter}
\else
    \usepackage{fontspec}
    \setmainfont{Charter}
\fi

% Settings
\raggedright
\AtBeginEnvironment{adjustwidth}{\partopsep0pt}
\pagestyle{empty}
\setcounter{secnumdepth}{0}
\setlength{\parindent}{0pt}
\setlength{\topskip}{0pt}
\setlength{\columnsep}{0.15cm}
\pagenumbering{gobble}

\titleformat{\section}{\needspace{3\baselineskip}\bfseries\large}{}{0pt}{}[\vspace{1pt}\titlerule]
\titlespacing{\section}{-1pt}{0.16cm}{0.08cm}

\renewcommand\labelitemi{$\vcenter{\hbox{\small$\bullet$}}$}

\newenvironment{highlights}{
  \begin{itemize}[topsep=0.04cm, parsep=0.02cm, partopsep=0pt, itemsep=0pt, leftmargin=10pt]
}{
  \end{itemize}
}

\newenvironment{onecolentry}{
  \begin{adjustwidth}{0.00001cm}{0.00001cm}
}{
  \end{adjustwidth}
}

\newenvironment{twocolentry}[2][]{
  \onecolentry
  \def\secondColumn{#2}
  \setcolumnwidth{\fill,4.5cm}
  \begin{paracol}{2}
}{
  \switchcolumn \raggedleft \secondColumn
  \end{paracol}
  \endonecolentry
}

\newenvironment{header}{
  \setlength{\topsep}{0pt}\par\kern\topsep\linespread{1.25}
}{
  \par\kern\topsep
}

\begin{document}
  % Header
  \begin{header}
    \centering
    \fontsize{22pt}{22pt}\selectfont Alex Dils

    \vspace{3pt}

    \small
    \faMapMarker\ Berkeley, CA \,|\,
    \href{mailto:dils@berkeley.edu}{\faEnvelope\ dils@berkeley.edu} \,|\,
    \href{https://www.alex-dils.com}{\faGlobe\ alex-dils.com} \,|\,
    \href{https://scholar.google.com/citations?user=0Sz8VPoAAAAJ}{\faGraduationCap\ Scholar} \,|\,
    \href{https://www.linkedin.com/in/alex-dils/}{\faLinkedin\ LinkedIn} \,|\,
    \href{https://github.com/axel-slid}{\faGithub\ GitHub}
  \end{header}

  \vspace{-0.18cm}

  \section{Education}
  \begin{twocolentry}{May 2028}
    \textbf{University of California, Berkeley}
  \end{twocolentry}
  B.A.\ in \textbf{\textit{Computer Science}} | GPA: 3.7/4.0

  \begin{onecolentry}
    \small \textit{Coursework:} Data Structures \& Algorithms, Machine Structures, Discrete Mathematics \& Probability, Multivariable Calculus
  \end{onecolentry}
  \begin{onecolentry}
    \small \textit{Activities:} Machine Learning at Berkeley (ML@B), Internal \& Research Committees
  \end{onecolentry}

  \section{Technical Skills}
  \begin{onecolentry}
    \small \textbf{Languages:} Python, JavaScript, TypeScript, Java, C, SQL; \textbf{ML:} PyTorch, TensorFlow, scikit-learn, Hugging Face, OpenCV, SimpleITK; \textbf{Web/Data:} FastAPI, Node.js, React, PostgreSQL, Supabase, PostGIS; \textbf{Tools:} Docker, Slurm, Git, Vercel
  \end{onecolentry}

  \section{Experience}

  

  \begin{twocolentry}{Jun 2026 -- Aug 2026}
    \textbf{\href{https://med.stanford.edu/canarycenter/canarycrest.html}{Stanford Canary CREST Research Scholar}}, Stanford Canary Center / Gevaert Lab
  \end{twocolentry}
  \begin{onecolentry}
    \begin{highlights}
      \item Developed and deployed volume estimation on cNF Vision, a clinical tool deployed to 50+ clinicians and patients.
      \item Improved lesion-volume estimation \textbf{$R^2$ from 0.78 to 0.85} on 40 held-out real images by building a physics-guided generative AI pipeline producing \textbf{42,305 synthetic 3D lesions} and \textbf{279 longitudinal sequences}.
      \item Reduced average depth-estimation RMSE by \textbf{72.9\%} (47.30 to 12.83 cm) on scans with simulated lesions by fine-tuning four depth models on synthetic 3D cNF data.
      \item Built a clinical tool for viewing and editing PET and CT scan annotations using JavaScript and Electron.
    \end{highlights}
  \end{onecolentry}
  
  \begin{twocolentry}{Feb 2026 -- Jun 2026}
    \textbf{Machine Learning Engineer},
    Logitech
  \end{twocolentry}
  \begin{onecolentry}
    \begin{highlights}
      \item Built a \textbf{VLM evaluation harness} for meeting-room readiness and participant classification on Logitech's Rally Board 65, measuring task accuracy and inference latency across 55 models.
      \item Ran model and prompt ablations and identified Qwen3.5 4B as the leading deployment candidate for the target hardware and task.
    \end{highlights}
  \end{onecolentry}
  
  \begin{twocolentry}{Jun 2022 -- Jun 2026}
    \textbf{Research Intern}, Stanford Center for Biomedical Informatics Research
  \end{twocolentry}
  \begin{onecolentry}
    \begin{highlights}
      \item Developed domain-robust skin-lesion segmentation methods and finite-element lung-motion augmentations that \textbf{improved SoTA tumor-segmentation Dice by >0.05}.
      \item Drafted sections of a Digital Twin manuscript and its journal submission cover letter; the research was \textbf{featured on Stanford Medicine's homepage}.
      \item Built a deployment tool for a model that generates corresponding spot-compression mammograms from CC and MLO views, making it deployable to clinicians and patients.
    \end{highlights}
  \end{onecolentry}

  % Projects
  \section{Projects}

  \begin{twocolentry}{2026}
    \textbf{Openleaf} \,|\, \href{https://www.alex-dils.com/openleaf}{alex-dils.com/openleaf}
  \end{twocolentry}
  \begin{onecolentry}
    \begin{highlights}
      \item Built a local macOS research workspace primarily in JavaScript, unifying LaTeX, PDF, Python, GitHub, SSH, and coding agents.
    \end{highlights}
  \end{onecolentry}

  \begin{twocolentry}{2026}
    \textbf{BsCode} \,|\, \href{https://www.alex-dils.com/bscode}{alex-dils.com/bscode}
  \end{twocolentry}
  \begin{onecolentry}
    \begin{highlights}
      \item Built a desktop primarily in JavaScript for coordinating Codex, Claude, and shell agents across local and remote workspaces.
    \end{highlights}
  \end{onecolentry}

  \begin{twocolentry}{2026}
    \textbf{LanternTrace Explorer} \,|\, \href{https://www.alex-dils.com/lanterntrace}{alex-dils.com/lanterntrace}
  \end{twocolentry}
  \begin{onecolentry}
    \begin{highlights}
      \item Built a physics-based explorer to track and forecast spotted lanternfly spread across the U.S. East Coast.
    \end{highlights}
  \end{onecolentry}

  % Publications
  \section{\texorpdfstring{Publications \hfill {\normalfont\itshape\footnotesize * indicates first authorship.}}{Publications}}

  \begin{onecolentry}
    \begin{highlights}
      \item \small \href{https://openreview.net/forum?id=JI0ZzOPW9x}{\textbf{*Synthesizing Longitudinal Cutaneous Neurofibroma Imaging for Digital-Twin Burden Quantification}}. MICCAI MWM 2026.\\
      Christoph Sadée, \textbf{Alex Dils}, Krish Sangani, Lillian Rubino, Alexander J Ferenchick, Varun Wadhwa, et al.
      \item \small \href{https://nhsjs.com/2025/eye-for-an-eye-a-deep-learning-and-analytical-method-to-spatializing-stereoscopic-images/\#google_vignette}{\textbf{*Eye For An Eye: A Deep-Learning and Analytical Method to Spatializing Stereoscopic Images}}. National High School Journal of Science, 2025.\\
      Jake Yoshinaka, \textbf{Alex Dils}, and Leo McDonnell.
      \item \small \textbf{*Addressing Bias and Confounders in AI-based Image Diagnosis --- A Study of 117,610 Skin Lesions}. Manuscript under revision at \textit{Cell}, 2025.\\
      Christoph Sadée, Serena Zhang, \textbf{Alex Dils}, Jared Bitz, Stefano Testa, Zhuo Ran Cai, et al.
      \item \small \href{https://scientificcomputing.github.io/fenics2024/integrating-mechanistic}{\textbf{Integrating Mechanistic Knowledge into Deep Learning for Improved Cancer Detection}}. FEniCS Conference Proceedings, 2024.\\
      Christoph Sadée, Katherine Hartmann, \textbf{Alex Dils}, Ianto Xi, Francisco Carrillo-Perez, et al.
    \end{highlights}
  \end{onecolentry}

\end{document}
